\documentclass[conference]{IEEEtran}
\usepackage{graphicx}
\usepackage{booktabs}
\usepackage{amsmath}
\usepackage{amssymb}
\usepackage{hyperref}
\usepackage{cite}
\usepackage{xcolor}
\usepackage{siunitx}

\graphicspath{{../results/figs/}}

\title{Adversarial Image Cloaking for Protection Against Face-Swap Attacks}

\author{
\IEEEauthorblockN{Jacob Antony Jeejo}
\IEEEauthorblockA{\textit{Indian Institute of Information Technology Kottayam}}
\and
\IEEEauthorblockN{Team Member 2}
\IEEEauthorblockA{\textit{Indian Institute of Information Technology Kottayam}}
\and
\IEEEauthorblockN{Team Member 3}
\IEEEauthorblockA{\textit{Indian Institute of Information Technology Kottayam}}
}

\begin{document}

\maketitle

\begin{abstract}
%TODO: Write abstract
This paper proposes a proactive defense mechanism against face-swap deepfakes by imperceptibly altering the source images to degrade the output quality of malicious models.
\end{abstract}

\section{Introduction}
%TODO: Write introduction
With the rapid advancements in generative AI, face-swap attacks pose significant privacy and security threats. We introduce a cloaking method based on adversarial perturbations to protect users from unauthorized manipulation.

\section{Related Work}
\subsection{Face Swapping}
%TODO: Describe face swapping methods
Recent models like SimSwap~\cite{chen2020simswap} and FaceShifter~\cite{li2019faceshifter} allow high-fidelity identity transfer.

\subsection{Proactive Defenses against Deepfakes}
%TODO: Proactive defenses
Methods such as Disrupting Deepfakes~\cite{ruiz2020disrupting} and Anti-Forgery~\cite{wang2022anti} have explored proactive manipulation against deepfakes, alongside CMUA-Watermark~\cite{huang2022cmua} for model-agnostic protection. Additional methods like PhotoGuard~\cite{salman2023photoguard} and Anti-DreamBooth~\cite{vanle2023anti} protect against broader AI editing.

\subsection{Anti-Face-Recognition Cloaking}
%TODO: Anti-FR cloaking
Fawkes~\cite{shan2020fawkes} and LowKey~\cite{cherepanova2021lowkey} are prominent works in cloaking identities against facial recognition systems, utilizing tools such as ArcFace~\cite{deng2019arcface} and InsightFace~\cite{insightface}.

\subsection{Purification and Adaptive Attacks}
%TODO: Purification
Diffusion-based purifications like DiffPure~\cite{nie2022diffpure} can potentially reduce adversarial effects, and gradient masking poses challenges as highlighted by BPDA~\cite{athalye2018obfuscated}.

\section{Threat Model}
%TODO: Threat model details
We assume the attacker has access to a face-swapping model and utilizes publicly available photos. The defender aims to release visually unaltered photos.

\section{Method}
\subsection{Identity-Embedding Attack}
%TODO: ID attack
To prevent face swapping, we disrupt the identity embeddings extracted by models such as FaceNet~\cite{schroff2015facenet} and Partial FC~\cite{an2021partial}.

\subsection{Differentiable Alignment}
%TODO: Alignment
We incorporate differentiable face alignment using methods based on standard pipelines like SCRFD~\cite{guo2022sample}.

\subsection{Ensemble and Targeted Objective}
%TODO: Ensemble objective
The PGD~\cite{madry2018towards} objective function is defined as:
\begin{equation}
\delta^* = \arg\min_{\|\delta\|_\infty \le \epsilon} \sum_k w_k \cos\big(E_k(A(x+\delta)), E_k(A(x))\big)
\end{equation}

\subsection{Expectation over Transformation}
%TODO: EOT
To improve robustness against image transformations, we use Expectation over Transformation (EOT)~\cite{athalye2018synthesizing} across resizing, blurring, and JPEG compression.

\subsection{Perceptual Constraint}
%TODO: Perceptual
To maintain image quality, we utilize the LPIPS~\cite{zhang2018unreasonable} metric alongside traditional bounds like FGSM~\cite{goodfellow2015explaining} and other properties~\cite{szegedy2014intriguing}.

\section{Experimental Setup}
%TODO: Experimental setup
Our experiments are conducted on datasets such as LFW~\cite{huang2007labeled}, VGGFace2~\cite{cao2018vggface2}, and FaceForensics++~\cite{rossler2019faceforensics}, commonly reviewed in deepfake surveys~\cite{tolosana2020deepfakes}.

\section{Results}
%TODO: Results discussion
Our method successfully degrades the output quality of face-swapping attempts while retaining high visual quality.
% \input{../results/tables/main.tex}

\section{Discussion and Limitations}
%TODO: Discussion
While effective against current models, the defense can potentially be bypassed by advanced purification.

\section{Ethical Considerations}
%TODO: Ethics
This work is intended to protect user privacy against malicious manipulations and does not aim to improve attacking capabilities.

\section{Conclusion}
%TODO: Conclusion
We presented a novel adversarial cloaking technique providing strong robustness against face-swap models.

\nocite{*}

\bibliographystyle{IEEEtran}
\bibliography{refs}

\end{document}
